\apisummary{
    Registers the arrival of a \ac{PE} at a synchronization point.
    This routine initiates a nonblocking synchronization operation for a 
    given \openshmem team and returns immediately without necessarily 
    completing the operation.
}

\begin{apidefinition}

\begin{C11synopsis}
int @\FuncDecl{shmem\_sync\_nb}@(shmem_team_t team, shmem_req_h *request);
\end{C11synopsis}

\begin{Csynopsis}
int @\FuncDecl{shmem\_team\_sync\_nb}@(shmem_team_t team, shmem_req_h *request);
\end{Csynopsis}

\begin{apiarguments}

\apiargument{IN}{team}{A valid \openshmem team handle to a team.}%
\apiargument{OUT}{request}{An opaque request handle identifying the synchronization
operation.}

\end{apiarguments}

\apidescription{
    \FUNC{shmem\_sync\_nb} is a collective nonblocking synchronization routine
    over an existing \openshmem team.

    {\bf Invocation and completion}: A call to the nonblocking sync routine
    initiates the operation and returns immediately without necessarily
    completing the operation. On success, an opaque request handle is created
    and returned. The completion of the operation can be observed after one or
    more calls to \FUNC{shmem\_req\_test} or a call to \FUNC{shmem\_req\_wait}.
    When the completion of the operation is observed, the request handle is
    deallocated and cannot be reused.

    {\bf Synchronization semantics}:
    The routine registers the arrival of a \ac{PE} at a synchronization point
    in the program. This is a fast mechanism for synchronizing all \acp{PE}
    that participate in this collective call. The routine ensures that all
    \acp{PE} in the specified team have called \FUNC{shmem\_sync\_nb} when
    the completion of the operation is observed.

    All \acp{PE} in the provided team must participate
    in the sync operation. If \VAR{team} compares equal to
    \LibConstRef{SHMEM\_TEAM\_INVALID} or is otherwise invalid, the behavior
    is undefined.

    Upon completion of the operation, the following is true for the local PE:
    \begin{itemize}
    \item All \acp{PE} in the team have called \FUNC{shmem\_sync\_nb}.
    \item Similar to the blocking \FUNC{shmem\_sync} routine, \FUNC{shmem\_sync\_nb}
    only ensures the completion and visibility of previously issued memory stores and
    does not ensure the completion of remote memory updates issued via OpenSHMEM routines.
    \end{itemize}
}

\apireturnvalues{
    Zero on successful local completion. Nonzero otherwise.
}

\apinotes{
    The \FUNC{shmem\_sync\_nb} routine can be used to portably ensure that
    memory access operations observe remote updates in the order enforced by the
    initiator \acp{PE}, provided that the initiator PE ensures completion of remote
    updates with a call to \FUNC{shmem\_quiet} prior to the call to the
    \FUNC{shmem\_sync\_nb} routine.

    Team handle error checking and integer return codes are currently undefined.
    Implementations may define these behaviors as needed, but programs should
    ensure portability by doing their own checks for invalid team handles and for
    \LibConstRef{SHMEM\_TEAM\_INVALID}.
}

\begin{apiexamples}

\apicexample
    {The following \Cstd[11] example is analogous to the \FUNC{shmem\_sync}
    example, but uses \FUNC{shmem\_sync\_nb} to overlap synchronization on two
    teams.}
    {./example_code/shmem_sync_nb_example.c}
    {}

\end{apiexamples}

\end{apidefinition}
